% ============================================================
%  国赛数学建模论文 LaTeX 模板
%  参考来源：肇庆学院校赛 C 题论文排版规范
%  编译方式：xelatex → xelatex（thebibliography 无需 bibtex）
% ============================================================
\documentclass[12pt,a4paper]{article}

% ============ 页面与字体 ============
\usepackage{geometry}
% Canonical defaults are recorded in config/paper_style.yaml; contest rules may override.
\geometry{a4paper,left=2.5cm,right=2.5cm,top=2.5cm,bottom=2.8cm,centering}
\usepackage{fontspec}
\setmainfont{SimSun}[BoldFont=SimHei,BoldFeatures={FakeBold=1.5}]
\setmonofont{SimSun}[AutoFakeBold=1.5]
\XeTeXlinebreaklocale "zh"
\XeTeXlinebreakskip = 0pt plus 2pt minus 0.5pt
\newfontfamily\heitifont{SimHei}[AutoFakeBold=1.5]
\newfontfamily\songfont{SimSun}[AutoFakeBold=1.5]
\newfontfamily\tnrfont{Times New Roman}

% ============ 数学与图形 ============
\usepackage{amsmath,amssymb,amsfonts}
\usepackage{graphicx,array}
\usepackage{xcolor}
\usepackage{listings}
\usepackage{float}
\usepackage{tabularx}
\usepackage{longtable}
\usepackage{booktabs}
\usepackage{makecell}
\usepackage{multirow}
\usepackage{calc}
\usepackage{tocloft}

% ============ 图表中文标签 ============
\renewcommand{\figurename}{图}
\renewcommand{\tablename}{表}
\renewcommand{\contentsname}{目录}
\renewcommand{\arraystretch}{1.25}
\renewcommand{\refname}{参考文献}

% ============ 标题格式（核心排版） ============
\usepackage{titlesec}
% 中文数字编号转换：1→一、2→二、...
\newcommand{\zhsectionnum}[1]{%
  \ifcase#1\or 一\or 二\or 三\or 四\or 五\or 六\or 七\or 八\or 九\or 十\else #1\fi}
% 编号格式
\renewcommand{\thesection}{\arabic{section}}
\renewcommand{\thesubsection}{\arabic{section}.\arabic{subsection}}
\renewcommand{\thesubsubsection}{\arabic{section}.\arabic{subsection}.\arabic{subsubsection}}
% 一级标题：居中、16pt黑体、"一、"前缀
\titleformat{\section}
  {\fontsize{15pt}{22pt}\selectfont\bfseries\heitifont\centering}
  {}{0em}{\zhsectionnum{\value{section}}、}
\titleformat{name=\section,numberless}
  {\large\bfseries\heitifont\centering}{}{0em}{}
\titlespacing{\section}{0pt}{18pt}{10pt}
% 二级标题：左对齐、正文大小黑体
\titleformat{\subsection}
  {\normalsize\bfseries\heitifont}{\thesubsection}{0.8em}{}
\titlespacing{\subsection}{0pt}{14pt}{8pt}
% 三级标题：左对齐、正文大小黑体
\titleformat{\subsubsection}
  {\normalsize\bfseries\heitifont}{\thesubsubsection}{0.8em}{}
\titlespacing{\subsubsection}{0pt}{10pt}{6pt}
\setcounter{secnumdepth}{3}

% ============ 段落与行距 ============
\setlength{\parindent}{2em}
\linespread{1.65}
\setlength{\parskip}{0pt}
\tolerance=800
\hbadness=4000
\frenchspacing
\widowpenalty=10000
\clubpenalty=10000
\raggedbottom
\pagestyle{plain} % no top-left running header
\setcounter{tocdepth}{2}

% 英文术语断词规则（消除表格 Overfull hbox）
\hyphenation{Ran-dom-For-est Grad-i-ent-Boost-ing Lo-gis-tic-Re-gres-sion}

% ============ 图表标题格式 ============
\makeatletter
\long\def\@makecaption#1#2{%
  \vskip 2pt%
  \begin{center}%
  \parbox{0.92\hsize}{\centering{\footnotesize\heitifont #1}\hspace{0.5em}{\footnotesize\heitifont #2}}%
  \end{center}%
  \vskip 2pt%
}
\makeatother

% ============ 引用上标（核心：正文 \cite{} 自动变右上角角标） ============
\makeatletter
\renewcommand{\@cite}[2]{\textsuperscript{[#1\if@tempswa , #2\fi]}}
\makeatother

% ============ 浮动体间距 ============
\setlength{\floatsep}{6pt plus 2pt minus 2pt}
\setlength{\textfloatsep}{7pt plus 2pt minus 2pt}
\setlength{\intextsep}{6pt plus 2pt minus 2pt}
\renewcommand{\topfraction}{0.9}
\renewcommand{\bottomfraction}{0.8}
\renewcommand{\textfraction}{0.08}
\renewcommand{\floatpagefraction}{0.75}
\AtBeginEnvironment{table}{\footnotesize}
\AtBeginEnvironment{longtable}{\footnotesize}
\setlength{\heavyrulewidth}{0.5pt}
\setlength{\lightrulewidth}{0.5pt}
\setlength{\aboverulesep}{0pt}
\setlength{\belowrulesep}{0pt}

% ============ 代码框格式（核心：附录代码排版） ============
\lstdefinestyle{modelcode}{
    language=Python,
    % Compact appendix box: denser than body text, but still readable in A4.
    basicstyle=\ttfamily\fontsize{8.5pt}{10.5pt}\selectfont,
    keywordstyle=\bfseries,
    commentstyle=\itshape\color{gray},
    numbers=left,
    numberstyle=\fontsize{6.5pt}{7.5pt}\selectfont\ttfamily\color{gray},
    stepnumber=1,
    numbersep=7pt,
    frame=single,
    framerule=0.35pt,
    rulecolor=\color{black!25},
    backgroundcolor=\color{gray!3},
    framesep=5pt,
    breaklines=true,
    breakatwhitespace=false,
    tabsize=4,
    showstringspaces=false,
    columns=fullflexible,
    keepspaces=true,
    xleftmargin=2.45em,
    xrightmargin=0.25em,
    framexleftmargin=2.1em
}
\lstdefinestyle{dgcupPython}{style=modelcode}
\renewcommand{\lstlistingname}{算法}
\lstnewenvironment{pycode}[1][]{\lstset{style=dgcupPython,#1}}{}

% ============ 表注命令 ============
\newcommand{\unirule}{\Xhline{0.8pt}}
\newcommand{\tabnote}[1]{%
  \par\vspace{0.3em}%
  \noindent
  \begin{minipage}{0.88\linewidth}%
  \raggedright\footnotesize\songfont #1%
  \end{minipage}%
  \par\vspace{0.5em}%
}

% ============ 超链接 ============
\usepackage[hidelinks]{hyperref}

% ============================================================
%  论文正文开始
% ============================================================
\newcommand{\PaperTitle}{在此填写论文标题}

\begin{document}

% ============ 首页：题目、摘要、关键词（无强制封面） ============
\thispagestyle{plain}
\begin{center}
  {\fontsize{16pt}{24pt}\selectfont\bfseries\heitifont \PaperTitle\par}
\end{center}
\begin{center}
  {\fontsize{14pt}{21pt}\selectfont\bfseries 摘\quad 要}
\end{center}\vspace{0.6cm}

% 摘要：问题本质与框架；逐题方法、关键结果、必要验证；适用范围。
在此填写摘要内容。第一段概述问题本质与总体框架；随后按题号写方法、最重要结果和必要验证。

\noindent{\bfseries 关键词：}关键词1；关键词2；关键词3；关键词4

\newpage
% ============ 目录 ============
\tableofcontents
\newpage

% ============ 正文 ============
\pagenumbering{arabic}\setcounter{page}{1}

% ---------- 一级标题示例：中文数字编号，居中 ----------
\section{问题重述}

% ---------- 二级标题示例：阿拉伯数字编号，左对齐 ----------
\subsection{问题背景}

在此填写问题背景。引用示例\cite{Si2020}。

\subsection{需要求解的问题}

在此列出需要求解的问题。

\section{问题分析}

\subsection{问题一分析}

在此分析问题一。可以引用多篇文献\cite{Hoerl1970,Bishop2006}。

\subsection{问题二分析}

在此分析问题二。

\section{模型假设}

所有假设如表~\ref{tab:assumptions}~所示。

\begin{center}\textbf{表1 模型假设}\end{center}

\begin{longtable}[]{@{}
  >{\centering\arraybackslash}p{(\linewidth - 4\tabcolsep) * \real{0.15}}
  >{\centering\arraybackslash}p{(\linewidth - 4\tabcolsep) * \real{0.35}}
  >{\centering\arraybackslash}p{(\linewidth - 4\tabcolsep) * \real{0.50}}@{}}
\toprule\noalign{}
\begin{minipage}[b]{\linewidth}\centering
\textbf{编号}
\end{minipage} & \begin{minipage}[b]{\linewidth}\centering
\textbf{假设}
\end{minipage} & \begin{minipage}[b]{\linewidth}\centering
\textbf{合理性论证}
\end{minipage} \\
\midrule\noalign{}
\endhead
\bottomrule\noalign{}
\endlastfoot
1 & 假设内容 & 论证内容 \\
2 & 假设内容 & 论证内容 \\
\end{longtable}
\label{tab:assumptions}

\tabnote{注：此处为表注示例，使用\texttt{\textbackslash tabnote}命令生成。}

\section{符号说明}

本文主要符号说明如表2所示。

% 仅当多个子问题确有共用状态、约束、数据或求解内核时，在此加入“统一模型层”。
\section{问题一}

\subsection{数据预处理与特征工程}

在此描述数据预处理方法。图~\ref{fig:example}~为示例图片。

\begin{figure}[htbp]\centering
  % \includegraphics[width=\textwidth]{figures/example.png}
  \fbox{\parbox{0.8\textwidth}{\centering\vspace{3cm}此处插入图片\vspace{3cm}}}
  \caption{示例图片标题}
  \label{fig:example}
\end{figure}

\subsection{模型建立}

在此建立数学模型。公式示例：
\begin{equation}
  \min_w \|Xw - y\|^2 + \alpha\|w\|^2
\end{equation}

\subsection{模型求解}

在此求解模型。

\subsection{结果分析}

在此分析结果。

\section{问题二}

% 按需继续添加子问题...

\section{模型检验与灵敏度分析}

在此给出约束审计、基线比较、独立复核、参数灵敏度或数值收敛中的必要证据，并解释其与结论的关系。

\section{模型的评价与推广}

\subsection{模型的优点}

（1）优点一。（2）优点二。

\subsection{模型的缺点}

（1）缺点一。（2）缺点二。

% ============ 参考文献 ============
% 【重要】按正文中首次出现顺序排列
% 格式：[编号] 作者. 标题[文献类型]. 出版信息, 年份.
\clearpage
\section{参考文献}

% the section title is already printed above; suppress the bibliography's
% second automatic heading while retaining its numbered entries.
\renewcommand{\refname}{}

\begin{thebibliography}{99}
\setlength{\itemsep}{2pt}
\setlength{\topsep}{0pt}
\setlength{\labelsep}{0.5em}
\bibitem{Si2020} 司守奎, 孙兆亮. 数学建模算法与应用(第3版)[M]. 北京: 国防工业出版社, 2020.
\bibitem{Hoerl1970} Hoerl A E, Kennard R W. Ridge Regression: Biased Estimation for Nonorthogonal Problems[J]. Technometrics, 1970, 12(1): 55-67.
\bibitem{Bishop2006} Bishop C M. Pattern Recognition and Machine Learning[M]. Springer, 2006.
% 继续添加参考文献，按首次引用顺序排列...
\end{thebibliography}

% ============ 附录 ============
\clearpage
\section{附录}

\subsection*{附录A\quad 数据预处理代码（preprocess.py 节选）}

以下为核心代码节选，完整代码见随附文件\texttt{src/preprocess.py}。

\begin{pycode}
import pandas as pd
import numpy as np
from pathlib import Path

def load_data(data_dir="data"):
    """加载原始数据"""
    data_dir = Path(data_dir)
    df = pd.read_csv(data_dir / "data.csv")
    return df
\end{pycode}

\subsection*{附录B\quad 模型训练代码（model.py 节选）}

以下为核心代码节选，完整代码见随附文件\texttt{src/model.py}。

\begin{pycode}
from sklearn.linear_model import Ridge

def train_model(X_train, y_train, alpha=1.0):
    """训练 Ridge 回归模型"""
    model = Ridge(alpha=alpha)
    model.fit(X_train, y_train)
    return model
\end{pycode}

\end{document}
